\chapter{Giới thiệu: Quá trình ngẫu nhiên \& Bước đi ngẫu nhiên\\{\normalsize\textit{Introduction: Stochastic Processes \& Random Walks}}}
\label{ch:intro}

% ============================================================================
\section{Quá trình ngẫu nhiên là gì? (What is a Stochastic Process?)}
\label{sec:what-is-sp}

\textbf{Câu chuyện:} Hãy tưởng tượng bạn đang đứng trên cầu nhìn xuống một chiếc lá trôi trên dòng sông.
Vị trí chiếc lá thay đổi theo thời gian, nhưng bạn không thể dự đoán chính xác nó sẽ ở đâu ---
dòng nước, gió, và xoáy nước đều tạo ra sự ngẫu nhiên.
Nếu ta ghi lại vị trí chiếc lá tại mỗi thời điểm, ta có một \term{quá trình ngẫu nhiên} (\en{stochastic process}).

\begin{dinhnghia}[(Quá trình ngẫu nhiên -- \en{Stochastic process})]
\label{def:stochastic-process}
Một \term{quá trình ngẫu nhiên} (\en{stochastic process}) là một họ các biến ngẫu nhiên
\[
  \{X_t : t \in T\}
\]
được định nghĩa trên cùng một không gian xác suất $(\Omega, \mathcal{F}, \PP)$,
trong đó $T$ là \term{tập chỉ số thời gian} (\en{index set}) và mỗi $X_t$ nhận giá trị trong
một \term{không gian trạng thái} (\en{state space}) $S$.
\end{dinhnghia}

\begin{tomtat}
Nói đơn giản: một quá trình ngẫu nhiên là một ``đại lượng thay đổi theo thời gian một cách ngẫu nhiên''.
Tại mỗi thời điểm $t$, ta có một biến ngẫu nhiên $X_t$ --- và tập hợp tất cả chúng tạo thành quá trình.
\end{tomtat}

\textbf{Hai loại thời gian:}
\begin{itemize}[nosep]
  \item \textbf{Thời gian rời rạc} (\en{discrete-time}): $T = \{0, 1, 2, \ldots\} = \N$.
    Ví dụ: giá cổ phiếu đóng cửa mỗi ngày, số khách hàng mỗi giờ.
  \item \textbf{Thời gian liên tục} (\en{continuous-time}): $T = [0, +\infty)$.
    Ví dụ: nhiệt độ tại một địa điểm, mực nước sông theo thời gian thực.
\end{itemize}

\textbf{Ba loại không gian trạng thái:}
\begin{itemize}[nosep]
  \item \textbf{Hữu hạn} (\en{finite}): $S = \{1, 2, \ldots, N\}$.
    Ví dụ: thời tiết $\in \{\text{nắng, mưa, mây}\}$.
  \item \textbf{Đếm được} (\en{countable}): $S = \Z$ hoặc $S = \N$.
    Ví dụ: số khách trong hàng đợi, vị trí bước đi ngẫu nhiên.
  \item \textbf{Không đếm được} (\en{uncountable}): $S = \R$ hoặc $S = \R^d$.
    Ví dụ: giá cổ phiếu, nhiệt độ.
\end{itemize}

\begin{dinhnghia}[(Đường mẫu -- \en{Sample path})]
\label{def:sample-path}
Cho một kết quả cố định $\omega \in \Omega$, ánh xạ
\[
  t \mapsto X_t(\omega)
\]
được gọi là một \term{đường mẫu} (\en{sample path}) hay \term{quỹ đạo} (\en{trajectory}) của quá trình.
Mỗi đường mẫu là một ``kịch bản'' cụ thể --- một lần ``chạy'' của quá trình.
\end{dinhnghia}

\begin{example}[Các ví dụ thực tế]
\label{ex:real-world-sp}
\begin{enumerate}[nosep]
  \item \textbf{Giá cổ phiếu:} $X_n$ = giá đóng cửa ngày thứ $n$. Thời gian rời rạc, không gian trạng thái liên tục.
  \item \textbf{Thời tiết:} $X_n \in \{\text{nắng, mưa}\}$. Thời gian rời rạc, không gian trạng thái hữu hạn.
  \item \textbf{Hàng đợi:} $X_t$ = số khách hàng trong hệ thống tại thời điểm $t$. Thời gian liên tục, không gian trạng thái đếm được.
  \item \textbf{DNA:} Dãy nucleotide $X_1, X_2, \ldots \in \{A, T, G, C\}$. Chỉ số là vị trí, không gian trạng thái hữu hạn.
\end{enumerate}
\end{example}

% ============================================================================
\section{Hai họ lớn: Markov và Martingale (Two Big Families)}
\label{sec:two-families}

Trong thế giới các quá trình ngẫu nhiên, hai họ quan trọng nhất là:

\textbf{1. Quá trình Markov} (\en{Markov process}): Tương lai chỉ phụ thuộc hiện tại, không phụ thuộc quá khứ.
\[
  \PP(X_{n+1} = j \mid X_n = i, X_{n-1}, \ldots, X_0) = \PP(X_{n+1} = j \mid X_n = i).
\]
Đây là tính chất ``không nhớ'' (\en{memoryless}). Khi không gian trạng thái là rời rạc,
ta gọi đây là \term{chuỗi Markov} (\en{Markov chain}) --- chủ đề chính của tài liệu này.
Ta sẽ nghiên cứu chi tiết từ Chương~1.

\textbf{2. Martingale}: Quá trình ``trò chơi công bằng'' --- kỳ vọng của giá trị tương lai,
khi biết quá khứ, đúng bằng giá trị hiện tại:
\[
  \E[X_{n+1} \mid X_n, X_{n-1}, \ldots, X_0] = X_n.
\]
Ý nghĩa: không có lợi thế hay bất lợi trong dài hạn. Ví dụ: tổng tiền thắng/thua trong một trò chơi công bằng.

\begin{tomtat}
\begin{center}
\begin{tabular}{lll}
\textbf{Họ} & \textbf{Ý tưởng chính} & \textbf{Ứng dụng} \\
\hline
Markov & ``Không nhớ'' & Chuỗi DNA, PageRank, hàng đợi \\
Martingale & ``Trò chơi công bằng'' & Tài chính, lý thuyết cờ bạc \\
\end{tabular}
\end{center}
Tài liệu này tập trung vào \textbf{chuỗi Markov}. Martingale chỉ xuất hiện ngắn gọn khi cần.
\end{tomtat}

% ============================================================================
\section{Bước đi ngẫu nhiên (Random Walks)}
\label{sec:random-walks}

Bước đi ngẫu nhiên là ví dụ đơn giản nhất --- và quan trọng nhất --- của cả quá trình Markov lẫn martingale.
Nó là ``Hello World'' của quá trình ngẫu nhiên.

\begin{dinhnghia}[(Bước đi ngẫu nhiên -- \en{Random walk})]
\label{def:random-walk}
Cho $(X_k)_{k \geq 1}$ là dãy các biến ngẫu nhiên \term{độc lập và cùng phân phối} (\en{i.i.d.}).
\term{Bước đi ngẫu nhiên} (\en{random walk}) được định nghĩa bởi $S_0 = 0$ và
\begin{equation}\label{eq:rw-def}
  S_n = X_1 + X_2 + \cdots + X_n, \qquad n \geq 1.
\end{equation}
\end{dinhnghia}

\begin{dinhnghia}[(Bước đi Bernoulli -- \en{Bernoulli random walk})]
\label{def:bernoulli-walk}
Khi mỗi bước $X_k$ chỉ nhận hai giá trị $\{+1, -1\}$ với
\[
  \PP(X_k = +1) = p, \qquad \PP(X_k = -1) = q = 1 - p,
\]
ta gọi $(S_n)_{n \geq 0}$ là \term{bước đi Bernoulli} (\en{Bernoulli random walk}).

Khi $p = q = 1/2$, ta gọi đây là bước đi \term{đối xứng} (\en{symmetric random walk}).
\end{dinhnghia}

\begin{example}[Người say rượu -- \en{Drunkard's walk}]
\label{ex:drunkard}
Một người say rượu đứng tại gốc tọa độ trên một con đường thẳng.
Mỗi bước, anh ta bước sang phải ($+1$) hoặc sang trái ($-1$) với xác suất bằng nhau $p = q = 1/2$.
Sau $n$ bước, vị trí của anh ta là $S_n$.

Câu hỏi tự nhiên: anh ta có quay lại gốc không? Đi xa bao nhiêu? Ta sẽ trả lời trong các mục tiếp theo.
\end{example}

% --- 3.1 Mean and Variance ---
\subsection{Trung bình và phương sai (Mean and Variance)}
\label{subsec:mean-var}

Ta tính kỳ vọng và phương sai của $S_n$ cho bước đi Bernoulli.

Kỳ vọng của mỗi bước:
\[
  \E[X_k] = (+1) \cdot p + (-1) \cdot q = p - q.
\]

Do đó, bằng tính tuyến tính của kỳ vọng:
\begin{equation}\label{eq:mean-Sn}
  \E[S_n] = n\,\E[X_1] = n(p - q).
\end{equation}

Phương sai của mỗi bước:
\begin{align*}
  \operatorname{Var}[X_k] &= \E[X_k^2] - (\E[X_k])^2 \\
  &= 1 \cdot q + 1 \cdot p - (p-q)^2 \\
  &= 1 - (2p-1)^2 = 4pq = 4p(1-p).
\end{align*}

Do các $X_k$ độc lập:
\begin{equation}\label{eq:var-Sn}
  \operatorname{Var}[S_n] = n\,\operatorname{Var}[X_1] = 4npq.
\end{equation}

\begin{congthuc}
\textbf{Trung bình và phương sai của bước đi Bernoulli:}
\[
  \E[S_n] = n(p - q), \qquad \operatorname{Var}[S_n] = 4npq.
\]
\begin{itemize}[nosep]
  \item Khi $p > q$: bước đi ``lệch phải'' (\en{positive drift}), $\E[S_n] > 0$.
  \item Khi $p < q$: bước đi ``lệch trái'' (\en{negative drift}), $\E[S_n] < 0$.
  \item Khi $p = q = 1/2$: bước đi đối xứng, $\E[S_n] = 0$, $\operatorname{Var}[S_n] = n$.
    Trung bình bằng $0$ nhưng phương sai tăng tuyến tính --- bước đi ``lang thang'' ngày càng xa gốc.
\end{itemize}
\end{congthuc}

% --- 3.2 Distribution ---
\subsection{Phân phối (Distribution)}
\label{subsec:distribution}

Ta muốn tính $\PP(S_n = k)$ cho bước đi Bernoulli.

Nhận xét quan trọng: $S_n$ chỉ có thể nhận các giá trị cùng tính chẵn lẻ với $n$.
Cụ thể, $S_{2n}$ chỉ có thể nhận giá trị chẵn, và $S_{2n+1}$ chỉ nhận giá trị lẻ.

Nếu $S_{2n} = 2k$, tức là trong $2n$ bước có $l = n+k$ bước đi lên và $2n - l = n - k$ bước đi xuống
(vì $l - (2n-l) = 2l - 2n = 2k$).

Xác suất của một đường đi cụ thể có $n+k$ bước lên và $n-k$ bước xuống là $p^{n+k}q^{n-k}$.
Số cách chọn vị trí cho $n+k$ bước lên trong $2n$ bước là $\binom{2n}{n+k}$.

\begin{dinhly}[(Phân phối của $S_{2n}$)]
\label{thm:dist-S2n}
\begin{equation}\label{eq:dist-S2n}
  \PP(S_{2n} = 2k) = \binom{2n}{n+k} p^{n+k}\,q^{n-k}, \qquad -n \leq k \leq n.
\end{equation}
\end{dinhly}

Đặc biệt, xác suất quay lại gốc tại thời điểm $2n$ là:
\begin{equation}\label{eq:return-origin}
  \PP(S_{2n} = 0) = \binom{2n}{n} p^n\,q^n = \binom{2n}{n}(pq)^n.
\end{equation}

\begin{example}[Bước đi đối xứng $p = q = 1/2$]
\label{ex:symmetric-dist}
Khi $p = q = 1/2$:
\[
  \PP(S_{2n} = 0) = \binom{2n}{n}\frac{1}{2^{2n}} = \binom{2n}{n}\frac{1}{4^n}.
\]
Sử dụng \term{xấp xỉ Stirling} (\en{Stirling's approximation}) $n! \approx \sqrt{2\pi n}\,(n/e)^n$:
\begin{equation}\label{eq:stirling-return}
  \PP(S_{2n} = 0) \approx \frac{1}{\sqrt{\pi n}}, \qquad n \to \infty.
\end{equation}
Xác suất này tiến về $0$ khi $n \to \infty$, nhưng rất chậm (tỉ lệ $1/\sqrt{n}$).
\end{example}

% --- 3.3 First Returns to Zero ---
\subsection{Quay lại gốc lần đầu (First Returns to Zero)}
\label{subsec:first-return}

Đặt
\begin{equation}\label{eq:tau0}
  \tau_0 = \inf\{n \geq 1 : S_n = 0\}
\end{equation}
là \term{thời gian quay lại gốc lần đầu} (\en{first return time to zero}).

\begin{dinhly}[(Xác suất quay lại gốc -- Privault, Mệnh đề 3.11)]
\label{thm:first-return}
Xác suất mà bước đi quay lại gốc trong thời gian hữu hạn là:
\begin{equation}\label{eq:prob-return}
  \PP(\tau_0 < \infty) = 2\min(p, q).
\end{equation}
Đặc biệt:
\begin{itemize}[nosep]
  \item Khi $p = q = 1/2$ (đối xứng): $\PP(\tau_0 < \infty) = 1$.
    Bước đi \textbf{chắc chắn} quay lại gốc.
  \item Khi $p \neq q$: $\PP(\tau_0 < \infty) < 1$.
    Có xác suất dương rằng bước đi \textbf{không bao giờ} quay lại gốc.
\end{itemize}
\end{dinhly}

\begin{tomtat}
Đây là một kết quả bất ngờ: bước đi đối xứng luôn quay lại gốc (xác suất 1),
nhưng \textbf{thời gian chờ trung bình} là vô hạn!
\[
  p = q = 1/2 \implies \PP(\tau_0 < \infty) = 1, \quad \text{nhưng} \quad \E[\tau_0] = +\infty.
\]
Kết quả này liên quan trực tiếp đến \textbf{tính hồi quy} (\en{recurrence}) của bước đi ngẫu nhiên 1 chiều,
mà ta sẽ gặp lại dưới dạng Định lý Polya trong Chương~3.
\end{tomtat}

% ============================================================================
\section{Từ bước đi ngẫu nhiên đến chuỗi Markov (From Random Walks to Markov Chains)}
\label{sec:rw-to-mc}

Bước đi ngẫu nhiên là một chuỗi Markov! Ta kiểm tra tính chất Markov:
do các $X_k$ độc lập, giá trị $S_{n+1} = S_n + X_{n+1}$ chỉ phụ thuộc vào $S_n$ và $X_{n+1}$,
không phụ thuộc vào $S_0, S_1, \ldots, S_{n-1}$:
\[
  \PP(S_{n+1} = j \mid S_n = i, S_{n-1}, \ldots, S_0) = \PP(S_{n+1} = j \mid S_n = i).
\]

Tuy nhiên, chuỗi Markov tổng quát hơn nhiều: xác suất chuyển $\PP(X_{n+1} = j \mid X_n = i)$
có thể là \textbf{bất kỳ giá trị nào} trong $[0,1]$ (miễn là tổng hàng bằng $1$),
không nhất thiết phải đối xứng hay đồng nhất.

\begin{congthuc}
\textbf{Bản đồ các chương:}
\begin{center}
\begin{tabular}{cl}
\textbf{Chương} & \textbf{Nội dung} \\
\hline
0 (đây) & Quá trình ngẫu nhiên, bước đi ngẫu nhiên \\
1 & Chuỗi Markov thời gian rời rạc, ma trận chuyển \\
2 & Phân tích bước đầu tiên: xác suất chạm, thời gian chạm \\
3 & Phân loại trạng thái: hồi quy, tạm thời, tuần hoàn \\
4 & Phân phối dừng và giới hạn \\
5 & Chuỗi Markov thời gian liên tục \\
A & Phụ lục: nền tảng xác suất \\
\end{tabular}
\end{center}
\end{congthuc}

\textbf{Câu chuyện cầu nối:} Hãy tưởng tượng một con kiến trên bàn cờ $3 \times 3$.
Mỗi bước, con kiến chọn ngẫu nhiên một ô kề cạnh để di chuyển.
Nếu bàn cờ chỉ có 1 chiều (một đường thẳng), con kiến thực hiện \emph{bước đi ngẫu nhiên}.
Nhưng trên bàn cờ 2 chiều, xác suất chuyển phụ thuộc vào vị trí
(ô góc có 2 hàng xóm, ô cạnh có 3, ô giữa có 4) --- đây là chuỗi Markov tổng quát.
Chương~1 sẽ khám phá thế giới này.

% ============================================================================
\section{Bài tập (Exercises)}
\label{sec:exercises-intro}

\begin{exercise}[Đường mẫu]
\label{ex:sample-paths}
Xét bước đi Bernoulli đối xứng ($p = q = 1/2$) với $n = 6$ bước.
\begin{enumerate}[nosep]
  \item Có bao nhiêu đường mẫu có thể? (Gợi ý: mỗi bước có 2 lựa chọn.)
  \item Liệt kê tất cả các đường mẫu kết thúc tại $S_6 = 2$.
  \item Tính $\PP(S_6 = 2)$ bằng hai cách: đếm trực tiếp và dùng công thức~\eqref{eq:dist-S2n}.
\end{enumerate}
\end{exercise}

\begin{exercise}[Phân phối $S_4$]
\label{ex:dist-S4}
Xét bước đi Bernoulli $(S_n)_{n \in \N}$ với $S_0 = 0$, xác suất tiến $p$ và xác suất lùi $1-p$.
\begin{enumerate}[nosep]
  \item Liệt kê tất cả các đường đi có thể dẫn đến $S_4 = 2$.
  \item Chứng minh rằng $\PP(S_4 = 2) = \binom{4}{3}p^3(1-p)$.
  \item Chứng minh công thức tổng quát: nếu $n+k$ chẵn và $|k| \leq n$,
    \[
      \PP(S_n = k) = \binom{n}{(n+k)/2} p^{(n+k)/2}\,(1-p)^{(n-k)/2}.
    \]
\end{enumerate}
\en{(Adapted from Privault, Exercise 3.1)}
\end{exercise}

\begin{exercise}[Trung bình và phương sai]
\label{ex:mean-var-calc}
Xét bước đi ngẫu nhiên $S_n = X_1 + \cdots + X_n$ với $\PP(X_k = +1) = 0.6$, $\PP(X_k = -1) = 0.4$.
\begin{enumerate}[nosep]
  \item Tính $\E[S_{100}]$ và $\operatorname{Var}[S_{100}]$.
  \item Sử dụng bất đẳng thức Chebyshev, tìm cận trên cho $\PP(|S_{100} - 20| \geq 30)$.
  \item Giải thích trực quan tại sao phương sai tăng theo $n$.
\end{enumerate}
\end{exercise}

\begin{exercise}[Quay lại gốc]
\label{ex:return-to-zero}
Cho bước đi Bernoulli đối xứng ($p = q = 1/2$).
\begin{enumerate}[nosep]
  \item Tính $\PP(S_2 = 0)$, $\PP(S_4 = 0)$, $\PP(S_6 = 0)$ bằng công thức~\eqref{eq:return-origin}.
  \item Kiểm tra rằng các giá trị trên phù hợp với xấp xỉ Stirling~\eqref{eq:stirling-return}.
  \item Giải thích tại sao $\PP(S_{2n+1} = 0) = 0$ cho mọi $n$.
\end{enumerate}
\end{exercise}

\begin{exercise}[Tính chất Markov của bước đi ngẫu nhiên]
\label{ex:rw-markov}
Cho bước đi Bernoulli $S_n = X_1 + \cdots + X_n$ với các $X_k$ i.i.d.
\begin{enumerate}[nosep]
  \item Chứng minh rằng $(S_n)_{n \geq 0}$ có tính chất Markov.

    \emph{Gợi ý:} Sử dụng tính độc lập của $X_{n+1}$ với $(X_1, \ldots, X_n)$.
  \item Viết ma trận chuyển $P$ cho bước đi Bernoulli trên không gian trạng thái $\Z$.
    Hàng $i$, cột $j$: $P_{i,j} = ?$
  \item Bước đi đối xứng ($p = 1/2$) có phải martingale không? Kiểm tra $\E[S_{n+1} \mid S_n] = S_n$.
\end{enumerate}
\end{exercise}
